\documentclass{article}
\usepackage[utf8]{inputenc}
\usepackage{amsmath, amssymb, amsthm}
\usepackage{geometry}
\usepackage{hyperref}
\geometry{a4paper, margin=2.5cm}
\title{HW5\_Report}
\date{\today}
\begin{document}
\maketitle
\documentclass{article}
\usepackage[utf8]{inputenc}
\usepackage{amsmath, amssymb, amsthm}
\usepackage{geometry}
\usepa‌​ckage{hyperref}
\geometry{a4paper, margin=1.5cm}
\title{HW5 Report}
\author{}
\date{\today}
\begin{document}
\maketitle
## Introduction

This report presents the analysis of the PET 212E Homework 5 data. The data includes measurements of porosity, grain density, and permeability at atmospheric conditions and at 800 psi.

## Porosity vs Grain Density Plot

The plot below shows the relationship between porosity and grain density at atmospheric conditions.

## Permeability vs Porosity Plot at 800 psi

The plot below shows the relationship between permeability and porosity at 800 psi.

## Conclusion

The plots provide valuable insights into the relationships between porosity, grain density, and permeability. These relationships are crucial in understanding the behavior of rocks in reservoir engineering applications.
\end{document}
\end{document}
